\documentclass[11pt]{article}

\usepackage[a4paper,margin=2.6cm]{geometry}
\usepackage[T1]{fontenc}
\usepackage[utf8]{inputenc}
\usepackage{lmodern}
\usepackage{microtype}
\usepackage{amsmath,amssymb}
\usepackage{graphicx}
\usepackage{booktabs}
\usepackage{caption}
\usepackage{xcolor}
\usepackage[colorlinks=true,linkcolor=blue!50!black,citecolor=blue!50!black,urlcolor=blue!50!black]{hyperref}

\captionsetup{font=small,labelfont=bf}

\newcommand{\occ}{\ensuremath{\mathrm{occ}}}
\newcommand{\dorm}{\ensuremath{\mathrm{dorm}}}

\title{Reinforcement Learning Teaches a Small Language Model to Use an
External Notepad Better than Its Own Context Window}

\author{Sten R\"udiger\\[2pt]
\small \url{https://github.com/sr-networks/clbench-memory-rl}}

\date{July 2026}

\begin{document}
\maketitle

\begin{abstract}
Agentic benchmarks of continual learning have reported a puzzling result: giving a
language model an external \emph{notepad} for carrying information across steps does
not reach the accuracy of plain in-context learning (ICL), in which the entire
interaction history simply remains in the prompt. We ask whether this deficit is a
property of the notepad interface or of the model's untrained \emph{policy} for using
it. On a long-horizon continual-monitoring task (30 sequential scans of a partially
observable radio band), we train Qwen3-1.7B with GRPO under a reward that pays
exclusively for recalling information that is no longer observable, and that is
orthogonal to single-step task skill by construction. We find that (i)~full-history
ICL is strongest only in the first few steps and then degrades as the transcript
grows, ending barely above a scripted no-memory upper bound; (ii)~the same base model
equipped with a notepad already surpasses ICL from step~6 onward \emph{without any
training}; and (iii)~reinforcement learning further lifts late-episode accuracy from
0.45 to 0.61 occupied-IoU---replicated in direction across four independent runs---%
while first-step accuracy, where the notepad is empty, remains unchanged
($0.216 \rightarrow 0.217$), confirming that the gain is memory use rather than task
skill. The reported notepad--ICL gap is therefore regime-dependent: over long
horizons a trained (and even an untrained) external memory dominates in-context
history.
\end{abstract}

\section{Introduction}
\label{sec:intro}

Language-model agents that act over long horizons must carry information across
steps. Two interface designs dominate practice. Under \emph{in-context learning}
(ICL), nothing is discarded: every past observation, action, and result stays in the
prompt, and the model re-reads its own history. Under an \emph{external
memory}~\cite{memgpt}, the context is windowed to the current step, and the agent
maintains a small, editable buffer---a \emph{notepad}---that is the only information
surviving from one step to the next.

A priori, the notepad appears preferable: it is a compressed, curated summary under
the agent's control, it keeps the prompt short, and it does not bury salient facts
under many turns of transcript. Yet the Continual-Learning Bench (CLBench), from
which our task is derived, reported the opposite ordering: notepad-equipped agents
did not reach the accuracy of full-history ICL. This is consistent with a simple
hypothesis that has, to our knowledge, not been isolated experimentally: pretrained
models have extensive implicit training in exploiting context, but essentially none
in \emph{maintaining} an external memory. The interface is not worse; the policy for
using it is untrained.

This paper tests that hypothesis with a controlled reinforcement-learning (RL)
study on a small model, Qwen3-1.7B~\cite{qwen3}. We pose one question:

\begin{quote}
\emph{Can RL train a small model to use an external notepad well enough to match or
exceed full-history ICL---without improving the model's single-step task
competence?}
\end{quote}

The italicized rider is essential. An RL run that improves a memory-dependent score
may have (a)~memorized task content into the weights, (b)~improved single-step
skill, or (c)~genuinely improved the carrying of information forward. Only (c) is
``training memory.'' Our experimental design excludes (a) and (b) by construction,
by guard, and by measurement (Section~\ref{sec:reward}).

Our contributions are: (1)~a memory-only reward for a long-horizon monitoring task,
red-teamed against blanket-reporting and weight-memorization exploits before
training; (2)~a four-condition comparison---scripted no-memory upper bound,
full-history ICL, untrained notepad, RL-trained notepad---aligned by step position
within the episode; and (3)~evidence that the published notepad--ICL gap reverses
over long horizons: ICL degrades as the transcript grows, consistent with known
long-context limitations~\cite{lostinthemiddle}, while the trained notepad accuracy
rises and holds.

\section{Task and Memory Interfaces}
\label{sec:task}

\paragraph{Blind spectrum monitoring.}
The task is derived from CLBench's blind-spectrum-monitoring family. An agent
monitors one fixed radio-frequency band containing 11--14 persistent transmitter
channels. It receives $n{=}30$ sequential scans of the same band. In any single
scan, only ${\sim}3$ transmitters are actively emitting, and the detector is noisy
(missed detections and false alarms). Dormant transmitters still \emph{occupy} the
band. After each scan the agent submits a report listing every frequency region it
believes to be persistently occupied (center frequency and bandwidth per region).

\paragraph{Score.}
The per-scan score is the occupied-spectrum intersection-over-union
(\occ-IoU) between the reported regions and the true persistent transmitter set,
\begin{equation}
\occ \;=\; \frac{\lvert R \cap T \rvert}{\lvert R \cup T \rvert} \in [0,1],
\end{equation}
where $R$ is the union of reported frequency intervals and $T$ the union of true
occupied intervals, both measured in MHz. Because only a minority of channels is
visible per scan, a high \occ{} requires reporting channels observed in
\emph{earlier} scans that are currently silent---i.e., memory.

\paragraph{No-memory upper bound.}
As the memoryless reference we use a scripted agent that reports exactly the
currently visible channels with perfect center frequencies and widths, replayed
through the real task engine on the same bands. This is not a weak baseline but an
\emph{upper bound} for any memoryless policy; it scores $\occ \approx 0.26$
(per-variant floors $a_0 \in \{0.2082, 0.2141, 0.2368\}$ enter the reward in
Section~\ref{sec:reward}).

\paragraph{Memory interfaces.}
We compare two designs holding the base model fixed:
\begin{itemize}
\item \textbf{ICL (full history).} Nothing is wiped; all past scans, reports, and
tool results remain in the prompt. No notepad exists.
\item \textbf{Notepad (windowed context).} The context contains only the current
scan. The agent has \texttt{notepad\_read}/\texttt{notepad\_write} tools and a
${\sim}4000$-character buffer it may overwrite each turn; this buffer is the only
information crossing scan boundaries. Maintaining it is the agent's own job.
\end{itemize}
The notepad system prompt teaches the \emph{procedure} (read the notepad, merge the
current peaks, write back the complete running set, report the full persistent
set); it does not contain any band content.

\section{A Memory-Only Reward}
\label{sec:reward}

RL on a memory-dependent score invites two families of confound: improving the
memory-free part of the task, and shortcut policies that fake coverage. The reward
is designed so that neither pays.

\paragraph{Currency.}
For scans $t = 2,\dots,n$ let the \emph{recallable set} $D_t$ be the channels that
are invisible at scan $t$ but were visible in some earlier scan. The per-scan
memory score is a Tversky overlap ($\alpha{=}\beta{=}1$) between the report and
$D_t$:
\begin{equation}
\mathrm{SCAN\_DORM}_t \;=\;
\frac{\mathrm{cov}(R_t, D_t)}
     {\mathrm{cov}(R_t, D_t) + \mathrm{FP}_t + \mathrm{FN}_t},
\end{equation}
where $\mathrm{cov}$ is the covered recallable area, $\mathrm{FN}_t$ the missed
recallable area, and---critically---$\mathrm{FP}_t$ the reported area on spectrum
the agent has \emph{never seen occupied}. Currently visible channels are excluded
from both numerator and penalties: reporting what is on screen is neither rewarded
nor punished. The currency is therefore orthogonal to single-step skill by
construction, and it is unearnable without information crossing scan boundaries.

\paragraph{Episode reward.}
One number per episode:
\begin{equation}
\label{eq:reward}
S = 3\left(\overline{\mathrm{dorm}}
 - p_{\mathrm{anchor}} - p_{\mathrm{carpet}} - p_{\mathrm{wmax}}
 - p_{\mathrm{complete}}\right),
\end{equation}
with $\overline{\mathrm{dorm}}$ the mean of $\mathrm{SCAN\_DORM}_t$ over the scans
actually played, and hinge penalties
\begin{align*}
p_{\mathrm{anchor}} &= \max(0,\, \mathrm{anchor} - a_0 - 0.10),
&& \mathrm{anchor} = \text{mean } \occ \text{ over scans 1--2},\\
p_{\mathrm{carpet}} &= 4.0 \cdot \max(0,\, \overline{\mathrm{rarea}} - 1.15),
&& \mathrm{rarea} = \text{reported area} / \text{true occupied area},\\
p_{\mathrm{wmax}} &= 1.0 \cdot \max(0,\, \overline{\mathrm{wmax}} - 1.5),
&& \mathrm{wmax} = \text{widest reported region} / \text{widest true channel},\\
p_{\mathrm{complete}} &= 1.5 \cdot \max(0,\, 30 - n_{\mathrm{played}} - 2)/30.
\end{align*}
An episode with zero completed scans (first-turn format failure) receives a flat
$S = -0.2$.

\paragraph{Adversarial audit.}
Each guard corresponds to an exploit simulated offline before training:
blanketing the band (``everything is occupied'') is caught by $p_{\mathrm{carpet}}$;
blanketing with one wide region under the area alarm by $p_{\mathrm{wmax}}$;
behaving honestly on scans 1--2 to satisfy the anchor and blanketing afterwards by
the per-scan area/width statistics; memorizing the fixed channel grid into the
weights (``gridbake'') by the definition of $\mathrm{FP}$ against never-seen space,
which an honest memory policy never paints; and truncating the episode after
banking coverage by $p_{\mathrm{complete}}$. Empirically all four penalties are
${\approx}0$ at epoch~0 in every run, so on honest play the reward is literally
$3\,\overline{\mathrm{dorm}}$---a memory-only currency. The anchor hinge
additionally pins near-memory-free scans to the scripted floor $a_0$, so a policy
that improves \emph{without} memory earns nothing and can be taxed. We deliberately
never reward the late-minus-early improvement itself, as rewarding a delta invites
sandbagging of early scans.

\paragraph{Weight-memorization detector.}
On scan~1 the notepad is empty, so any trained-versus-untrained gain there must be
band knowledge stored in the weights rather than memory use. Scan-1 \occ{} is
therefore measured for every run as a no-memory discriminator; the result
(Section~\ref{sec:results}) is flat to within $+0.001$.

\section{Experimental Setup}
\label{sec:setup}

\paragraph{Training.}
Base model Qwen3-1.7B~\cite{qwen3}; GRPO~\cite{grpo} with LoRA adapters on a
managed reinforcement fine-tuning service. Four replicate runs with identical
configuration and different seeds (job IDs \texttt{s8e07n53}, \texttt{yp8deoer},
\texttt{kym4znjc}, \texttt{wqjyy66p}): 24 training bands, learning rate
$2{\times}10^{-4}$, 5 epochs, 8 candidate rollouts per GRPO group, sampling
temperature 1.2, and a 16{,}384-token per-turn output cap. Rewards follow
Eq.~\eqref{eq:reward}. Per epoch and run, evaluation pools 192 rollouts; ``trained''
denotes the final epoch (ep4), ``untrained'' the same runs at ep0.

\paragraph{Conditions.}
The central comparison aligns four conditions by scan position within the episode:
\begin{center}
\small
\begin{tabular}{@{}llll@{}}
\toprule
condition & policy & memory & rollouts \\
\midrule
no-mem & scripted perfect memoryless agent & none (upper bound) & deterministic replay \\
ICL & untrained base, full history in prompt & in-context & 288 (job \texttt{g7dncu2c}, ep0) \\
notepad-untrained & untrained base with notepad tools & notepad & $4\times192$ (ep0, pooled) \\
notepad-trained & same four runs after RL & notepad & $4\times192$ (ep4, pooled) \\
\bottomrule
\end{tabular}
\end{center}
ICL-versus-notepad-untrained is a pure comparison of memory \emph{interface} (same
weights); notepad-untrained-versus-trained isolates the \emph{training} effect.
Scores are pooled into six bins of five consecutive scans.

\section{Results}
\label{sec:results}

\begin{figure}[t]
\centering
\includegraphics[width=\textwidth]{occ_by_scan_bin.png}
\caption{Occupied-IoU by scan position within the 30-scan episode, for the four
conditions. Full-history ICL (purple) peaks at scans 6--10 and then degrades as the
transcript grows, ending barely above the scripted no-memory upper bound (grey).
The untrained notepad (light blue) surpasses ICL from scan~6 onward; RL
(green; whiskers give the min--max of the four replicate runs) lifts late-episode
accuracy to ${\approx}0.61$.}
\label{fig:main}
\end{figure}

\begin{table}[t]
\centering
\caption{Occupied-IoU by scan-position bin (higher is better). Notepad columns pool
four replicate runs; the bracketed range is the min--max of the four per-run means
at ep4. $n$ is the number of trained-condition scan observations pooled per bin.
Episode-level bootstrap 95\% confidence intervals of the pooled bin means (2000
resamples of whole episodes) are no wider than $\pm 0.011$ for the notepad
conditions and $\pm 0.022$ for any condition; beyond the first bin the four
conditions' intervals are pairwise non-overlapping. Full CIs ship with the data
(Data availability).}
\label{tab:main}
\small
\begin{tabular}{@{}lccccc@{}}
\toprule
scans & no-mem & ICL & notepad ep0 & notepad ep4 & $n_{\mathrm{ep4}}$ \\
\midrule
1--5   & 0.257 & \textbf{0.348} & 0.332 & 0.394 \,[0.301--0.478] & 3840 \\
6--10  & 0.264 & 0.378 & 0.436 & \textbf{0.559} \,[0.421--0.656] & 3840 \\
11--15 & 0.268 & 0.371 & 0.431 & \textbf{0.581} \,[0.463--0.666] & 3825 \\
16--20 & 0.270 & 0.308 & 0.455 & \textbf{0.595} \,[0.501--0.662] & 3808 \\
21--25 & 0.268 & 0.294 & 0.454 & \textbf{0.615} \,[0.540--0.679] & 3737 \\
26--30 & 0.266 & 0.286 & 0.453 & \textbf{0.612} \,[0.553--0.678] & 3171 \\
\bottomrule
\end{tabular}
\end{table}

Figure~\ref{fig:main} and Table~\ref{tab:main} contain the main result.

\paragraph{ICL degrades over the episode.}
Full-history ICL leads only in the first bin (0.348 vs.\ 0.332 for the untrained
notepad), peaks at scans 6--10 (0.378), and then falls monotonically---0.371,
0.308, 0.294, 0.286---although strictly more information is present at every
scan. By the final bin it is barely above the no-memory bound (0.266), and its
late-half mean (scans 16--30) is 0.297. This is the long-context failure mode
at small scale: facts buried twenty-plus scans deep in a growing transcript
effectively stop being used, consistent with positional-use limitations of long
contexts~\cite{lostinthemiddle}.

\paragraph{The untrained notepad already beats ICL beyond scan 5.}
With identical weights, windowing the context and providing the notepad yields
${\approx}0.45$ \occ{} throughout the late half---$+0.16$ over ICL---before any
training. A short, curated summary pinned at a fixed position does not rot the
way raw history does. This alone reframes the benchmark finding that motivated the
study: the notepad interface is not inherently weaker; over long horizons it is
inherently more robust.

\paragraph{RL lifts the notepad to a decisive win.}
Training raises late-half \occ{} from 0.454 to 0.607 (pooled), a $+0.31$ margin
over ICL and roughly a doubling of the headroom above the no-memory floor. The
direction of the effect replicated in four of four runs (complete-episode mean
recallable-coverage $+0.041$, $+0.101$, $+0.181$, $+0.182$); the magnitude varies
by seed, which the whiskers in Figure~\ref{fig:main} report rather than hide.

\paragraph{The gain is memory, not skill.}
Scan-1 \occ{}, measured on an empty notepad, moves by $+0.0004$ to $+0.0011$
across the four runs (pooled $0.216 \rightarrow 0.217$): the trained model is not
better at the task when it has nothing to remember. Combined with the construction
of the currency (visible channels carry no gradient), the anchor pin, and the
silent guards, every improvement in Table~\ref{tab:main} is attributable to
carrying information forward.

\paragraph{Training ICL does not fix the collapse.}
For completeness we also trained an ICL arm under a dense per-scan accuracy reward
(job \texttt{g7dncu2c}, 4096-token cap). It learned to degrade somewhat less over
the long context (mean \occ{} $0.344 \rightarrow 0.363$), but its late-minus-early
\occ{} remained negative in every epoch: it never learned to accumulate. The
degradation of in-context memory over long horizons appears to be the binding
constraint at this model scale, not an artifact of the untrained comparison.

\section{Limitations}
\label{sec:limitations}

\emph{Survivorship in late bins.} At ep4, episode completion across the four runs
is 100\%, 97\%, 60\%, and 90\%; one run truncates episodes early. Late bins
therefore average only scans actually played ($n$ falls from 3840 to 3171), a
disclosed survivorship bias. \emph{ICL condition.} The ICL bar is the untrained
base model under a 4096-token per-turn cap; the trained-ICL check above mitigates
but does not remove the asymmetry. \emph{Scale and scope.} One base model
(Qwen3-1.7B), one optimizer (GRPO), one task family; learning rate, temperature,
and group size were not swept. \emph{Benchmark relation.} This is a
re-implementation of the CLBench task family as an RL environment, not the
benchmark's exact harness; the contribution is the mechanism and its
trainability, not a leaderboard number. \emph{Run selection.} The four reported
runs are the terminal configuration of a longer design sequence, not survivors
selected from it: a registry of every training job bought for this project (43,
including all failed reward designs and a failed replication of an earlier
configuration) is published with the data.

\section{Conclusion}
\label{sec:conclusion}

On a long-horizon continual-monitoring task, the reported accuracy advantage of
in-context history over an external notepad is a short-horizon phenomenon. As the
transcript grows, ICL degrades toward the no-memory bound, while a notepad-equipped
agent---even untrained---holds its accuracy, and reinforcement learning under a
memory-only, adversarially audited reward widens the gap to $+0.31$ \occ-IoU in the
late episode, replicated in direction across four seeds, with first-step accuracy
provably unchanged. For small models acting over long horizons, the practical
implication is direct: do not hoard history---train the model to write things down.

\subsection*{Data availability}
All numbers in Table~\ref{tab:main} and Figure~\ref{fig:main} are recomputable
from raw data. The repository at
\url{https://github.com/sr-networks/clbench-memory-rl} contains the raw
per-episode, per-scan traces of all four conditions (54{,}420 rows), a script
that regenerates the figure and the binned table---including the bootstrap
confidence intervals---from those traces alone, the executed reward code
verbatim, the scripted no-memory agent, the task description, and the complete
run registry. Training runs are identified by their training-job IDs throughout,
so each number is traceable to its source job.

\begin{thebibliography}{9}

\bibitem{lostinthemiddle}
N.~F. Liu, K.~Lin, J.~Hewitt, A.~Paranjape, M.~Bevilacqua, F.~Petroni, and
P.~Liang.
\newblock Lost in the middle: How language models use long contexts.
\newblock \emph{Transactions of the Association for Computational Linguistics},
12:157--173, 2024.

\bibitem{grpo}
Z.~Shao, P.~Wang, Q.~Zhu, R.~Xu, J.~Song, X.~Bi, H.~Zhang, M.~Zhang, Y.~K. Li,
Y.~Wu, and D.~Guo.
\newblock DeepSeekMath: Pushing the limits of mathematical reasoning in open
language models.
\newblock \emph{arXiv preprint arXiv:2402.03300}, 2024.

\bibitem{memgpt}
C.~Packer, S.~Wooders, K.~Lin, V.~Fang, S.~G. Patil, I.~Stoica, and
J.~E. Gonzalez.
\newblock MemGPT: Towards LLMs as operating systems.
\newblock \emph{arXiv preprint arXiv:2310.08560}, 2023.

\bibitem{qwen3}
Qwen Team.
\newblock Qwen3 technical report.
\newblock \emph{arXiv preprint arXiv:2505.09388}, 2025.

\end{thebibliography}

\end{document}
